\section{LEVEL 0 · Startklaar}

\begin{missie}
Missie: leer de taal van de statistiek spreken.
\end{missie}

\begin{doel}
\textbf{DOEL} \quad Na dit level kan je:

☐ de begrippen \textbf{populatie, steekproef, eenheid, variabele} en \textbf{waarneming} gebruiken;

☐ het \textbf{soort variabele} en het \textbf{soort data} bepalen;

☐ een \textbf{frequentietabel} opstellen en lezen (AF en RF).
\end{doel}

\subsection*{Theorie 0.1 · Basisbegrippen}

\begin{theorie}
\textbf{Statistiek} is data verzamelen, ordenen, voorstellen en interpreteren, om er besluiten uit te trekken.

\begin{itemize}
\item De \textbf{populatie} is de volledige groep waarover je een uitspraak wil doen. Aantal: $N$.
\item De \textbf{steekproef} is het deel van de populatie dat je effectief onderzoekt. Aantal: $n$.
\item Een \textbf{eenheid} is één persoon of voorwerp uit de groep.
\item De \textbf{variabele} is het kenmerk dat je onderzoekt.
\item Een \textbf{waarneming} (gegeven) is de waarde van de variabele bij één eenheid. Alle waarnemingen samen vormen de \textbf{dataset}.
\end{itemize}
\end{theorie}

\begin{voorbeeld}
\textbf{Voorbeeld.} Een kinderdagverblijf wil weten hoeveel uur per week de kinderen er verblijven. Er zijn 250 kinderen ingeschreven. Het team bevraagt de ouders van 40 kinderen.

\begin{itemize}
\item Populatie: alle 250 kinderen van het kinderdagverblijf ($N = 250$).
\item Steekproef: de 40 bevraagde kinderen ($n = 40$).
\item Eenheid: één kind.
\item Variabele: het aantal uur verblijf per week.
\item Waarneming: bijvoorbeeld ‘38 uur’.
\end{itemize}
\end{voorbeeld}

\subsection*{Theorie 0.2 · Soorten variabelen en soorten data}

\begin{theorie}
Een \textbf{kwalitatieve} variabele beschrijft een eigenschap met woorden of labels (data zijn \textbf{categorisch}). Een \textbf{kwantitatieve} variabele is een getal waarmee je zinvol kan rekenen (data zijn \textbf{numeriek}).

\breedtes{24/28/36}
\begin{tabular}{lll}
\hline
Soort variabele & Soort data & Voorbeeld \\
\hline
kwalitatief & categorisch niet-geordend & lievelingskleur, soort school (ASO, TSO, BSO, KSO) \\
kwalitatief & categorisch geordend & tevredenheid van ‘zeer ontevreden’ tot ‘zeer tevreden’ \\
kwantitatief & numeriek, \textbf{discreet} (je telt) & aantal broers en zussen \\
kwantitatief & numeriek, \textbf{continu} (je meet) & lengte, tijd, gewicht \\
\hline
\end{tabular}
\end{theorie}

\begin{letop}
\textbf{LET OP} \quad Een getal is niet altijd numerieke data. Een rugnummer, een postcode of een telefoonnummer is gewoon een \emph{label}: rekenen met zo’n getal (bijvoorbeeld een gemiddelde nemen) heeft geen enkele betekenis.
\end{letop}

\subsection*{Theorie 0.3 · De frequentietabel}

\begin{theorie}
Een \textbf{frequentietabel} geeft weer hoe vaak elke waarde voorkomt.

\begin{itemize}
\item De \textbf{absolute frequentie} (AF) is het aantal keer dat een waarde voorkomt. De som van alle AF is $n$.
\item De \textbf{relatieve frequentie} (RF) is het aandeel in de steekproef: $\;\text{RF}=\dfrac{\text{AF}}{n}\cdot 100\,\%$.
\end{itemize}
\end{theorie}

\begin{voorbeeld}
\textbf{Voorbeeld.} In een kleuterklas wordt genoteerd in welke hoek de 20 kleuters tijdens het vrij spel spelen.

\breedtes{30/18/18}
\begin{tabular}{lll}
\hline
speelhoek & AF & RF (\%) \\
\hline
bouwhoek & 6 & 30,0 \\
poppenhoek & 8 & 40,0 \\
tekenhoek & 4 & 20,0 \\
leeshoek & 2 & 10,0 \\
\textbf{totaal} & \textbf{20} & \textbf{100,0} \\
\hline
\end{tabular}

Berekening: $\dfrac{6}{20}\cdot 100\,\% = 30\,\%$. De variabele ‘speelhoek’ is kwalitatief (categorisch niet-geordend).
\end{voorbeeld}

\subsection*{Aan de slag}

\subsubsection*{Oefening 0.1 \quad ★☆☆ \quad Invullen: de basisbegrippen}

Vul de ontbrekende woorden in. \textbf{Woordenlijst:} populatie – steekproef – variabele – eenheid – waarneming.

Een \gat{populatie} is de volledige groep waarover je iets wil weten. Als je niet iedereen kan onderzoeken, bevraag je een deel van die groep: de \gat{steekproef}. Het kenmerk dat je onderzoekt, noem je de \gat{variabele}. Eén persoon uit de groep is een \gat{eenheid}. De waarde van de variabele bij die persoon is een \gat{waarneming}.

\subsubsection*{Oefening 0.2 \quad ★★☆ \quad Koppel: welk soort variabele?}

Vul de tabel aan. Gebruik de woorden uit theorie 0.2. Bij numerieke data geef je ook aan of ze discreet of continu zijn.

\breedtes{38/20/22/12}
\begin{tabular}{llll}
\hline
Variabele & kwalitatief of kwantitatief? & soort data & discreet of continu? \\
\hline
a) het aantal kinderen in een gezin & & & \\
b) de lievelingskleur van een kleuter & & & \\
c) de tevredenheid over de opvang (van ‘zeer ontevreden’ tot ‘zeer tevreden’) & & & \\
d) de lichaamstemperatuur van een baby & & & \\
e) het soort onderwijs (ASO, TSO, BSO, KSO) & & & \\
f) de wachttijd in minuten op de spoeddienst & & & \\
g) het rugnummer van een voetballer & & & \\
\hline
\end{tabular}

\begin{oplossing}
\textbf{Oplossing 0.2}

a) kwantitatief – numeriek – discreet \quad b) kwalitatief – categorisch niet-geordend \quad c) kwalitatief – categorisch geordend \quad d) kwantitatief – numeriek – continu \quad e) kwalitatief – categorisch niet-geordend \quad f) kwantitatief – numeriek – continu \quad g) kwalitatief – categorisch niet-geordend (een rugnummer is een label, geen hoeveelheid).
\end{oplossing}

\subsubsection*{Oefening 0.3 \quad ★★☆ \quad De frequentietabel invullen}

De 25 leerlingen van een klas worden gevraagd hoe ze naar school komen.

\breedtes{30/18/18}
\begin{tabular}{lll}
\hline
vervoermiddel & AF & RF (\%) \\
\hline
te voet & 5 & \\
fiets & 8 & \\
bus & 7 & \\
auto & 3 & \\
trein & 2 & \\
\textbf{totaal} & & \\
\hline
\end{tabular}

\textbf{a)} Vul de tabel aan. Rond af op 1 decimaal indien nodig.

\textbf{b)} Welk percentage van de klas komt met de fiets? Schrijf een antwoordzin. \lijn

\textbf{c)} Wat is hier de populatie? En wat is het soort variabele? \lijn

\begin{oplossing}
\textbf{Oplossing 0.3}

a) RF: 20,0 – 32,0 – 28,0 – 12,0 – 8,0; totaal AF = 25 en totaal RF = 100,0. Bijvoorbeeld fiets: $\frac{8}{25}\cdot 100\,\% = 32\,\%$.

b) 32 \% van de leerlingen van de klas komt met de fiets naar school.

c) Populatie: de 25 leerlingen van de klas (als je alleen over deze klas een uitspraak doet; zijn de 25 een deel van de hele school, dan zijn zij een steekproef). Variabele: ‘vervoermiddel’, kwalitatief – categorisch niet-geordend.
\end{oplossing}

\subsubsection*{Oefening 0.4 \quad ★★☆ \quad Context: wie, wat, hoeveel?}

Vul de tabel aan voor elke situatie.

\breedtes{30/16/16/16/22}
\begin{tabular}{lllll}
\hline
Situatie & populatie & steekproef & variabele & soort variabele \\
\hline
a) Een jeugdwerker wil weten hoeveel uur per dag 12-jarigen in Vlaanderen gamen. Hij vraagt het aan 150 12-jarigen uit vijf scholen. & & & & \\
b) Een woonzorgcentrum vraagt aan 30 van zijn 120 bewoners of ze ‘zeer tevreden’, ‘tevreden’, ‘ontevreden’ of ‘zeer ontevreden’ zijn over de maaltijden. & & & & \\
\hline
\end{tabular}

\begin{oplossing}
\textbf{Oplossing 0.4}

a) populatie: alle 12-jarigen in Vlaanderen; steekproef: de 150 bevraagde 12-jarigen; variabele: aantal uur gamen per dag; kwantitatief (continu).

b) populatie: alle 120 bewoners; steekproef: de 30 bevraagde bewoners; variabele: tevredenheid over de maaltijden; kwalitatief (categorisch geordend).
\end{oplossing}

\subsection*{CHECK · Level 0}

\textbf{1.} (2 p) Is ‘het aantal kinderen in een klas’ een kwalitatieve of een kwantitatieve variabele? Is ze discreet of continu? \lijn

\textbf{2.} (2 p) In een steekproef van 40 jongeren zeggen er 12 dat ze elke dag sporten. Bereken de relatieve frequentie in \% en schrijf een antwoordzin. \lijn

\textbf{3.} (2 p) Een school wil weten hoe tevreden de 600 leerlingen zijn over de speeltijd. Ze bevragen er 60. Noteer $N$ en $n$ en benoem de steekproef. \lijn

\textbf{4.} (2 p) Geef een voorbeeld van een variabele met categorisch geordende data. \lijn

\begin{doel}
\textbf{Score:} \_\_\_ / 8 \quad ☐ behaald (vanaf 6 / 8) \quad ☐ nog even oefenen
\end{doel}

\begin{oplossing}
\textbf{Oplossing CHECK 0}

1. Kwantitatief en discreet (je telt kinderen). 

2. $\text{RF} = \frac{12}{40}\cdot 100\,\% = 30\,\%$. Antwoordzin: 30 \% van de bevraagde jongeren sport elke dag.

3. $N = 600$ en $n = 60$; de steekproef bestaat uit de 60 bevraagde leerlingen.

4. Bijvoorbeeld tevredenheid (zeer ontevreden … zeer tevreden), een rangorde, een score als ‘laag – gemiddeld – hoog’.
\end{oplossing}
